\section{Interest: The Shape of Cost Growth Over Time}
\label{sec:interest}

\Cref{prop:steps} predicts the temporal signature of reconstruction effort and,
under a sound exhaustive policy, abstention exposure: plateaus punctuated by cliffs, with smooth
growth where a hazard-governed human path is the cheapest surviving source. Error
exposure may instead drift as discriminators degrade. This section makes the
qualified prediction concrete by computing
the expected-cost trajectory of one canonical deficit under a documented decay
schedule. The computation is \emph{model-derived and clearly labeled as such}: its
inputs are realistic and cited, its output is an illustration of structure, not a
field measurement.

\subsection{A Worked Decay Schedule}
\label{sec:interest-worked}

Consider the paradigmatic deficit of practice P1 (\cref{sec:practices}): a change
deployed with no intent link---the ``why'' was never durably attached. Four
reconstruction paths exist at creation, with survival forms from
\cref{sec:model-decay}: the CI log, whose narrative context often reveals intent
(cheap; step survival at a 90-day retention default~\cite{nist80092,ghdocs}); forge
context---branch names, adjacent commits, review discussion (moderate cost; event
survival until a forge migration, here placed at day 400, an instance of P7); the
author's memory (expensive, and increasingly so with elapsed time; hazard survival
with mean tenure $\sim$2.5 years~\cite{rigby2016}); and nothing (irrecoverable).

\begin{figure}[t]
\centering
% Interest trajectory (single column), data from ../data/interest.csv
\begin{tikzpicture}
\begin{axis}[
  edplot,
  width=0.92\columnwidth, height=5.1cm,
  xlabel={deficit age $\tau$ (days)},
  ylabel={expected reconstruction cost (log)},
  ymode=log, ymin=0.5, ymax=200,
  xmin=0, xmax=1100,
  legend pos=north west,
  axis y line*=left,
]
\addplot[serA, thick]
  table[x=day, y=exp_cost, col sep=comma] {data/interest.csv};
\addlegendentry{expected cost}
\coordinate (cliff1) at (axis cs:90,0.5);
\coordinate (cliff2) at (axis cs:400,0.5);
\draw[black!35, densely dashed] (axis cs:90,0.5) -- (axis cs:90,200);
\draw[black!35, densely dashed] (axis cs:400,0.5) -- (axis cs:400,200);
\node[lbl, anchor=south, rotate=90] at (axis cs:78,2.6) {CI retention cliff (90 d)};
\node[lbl, anchor=south, rotate=90] at (axis cs:388,2.6) {forge migration (400 d)};
\end{axis}
\begin{axis}[
  edplot,
  width=0.92\columnwidth, height=5.1cm,
  xmin=0, xmax=1100,
  ymin=0, ymax=1.02,
  axis y line*=right,
  axis x line=none,
  ylabel={$\Pr[\text{irrecoverable}]$},
  legend pos=south east,
]
\addplot[serB, thick, densely dashed]
  table[x=day, y=p_irrec, col sep=comma] {data/interest.csv};
\addlegendentry{$\Pr[\text{irrecoverable}]$}
\end{axis}
\end{tikzpicture}

\caption{Model-derived interest trajectory for a fixed intent-link deficit under
the decay schedule of \cref{sec:interest-worked} (computed from the accompanying
script; illustrative structure, not a field measurement). The expected evidence-
debt contribution for the illustrated fixed policy,
$B_{q,\rho}(\tau)=\mathbb{E}C_\rho+\lambda_q\Pr_\rho(\Err_q)
+\pi_q\Pr_\rho(\mathsf{Abst}_q)$ (solid, left axis, log scale), is flat while the cheapest path (CI log)
survives, jumps at its 90-day retention cliff to the forge-context cost, jumps
again at the day-400 migration to the interview path---whose expected cost then
grows smoothly as turnover hazard accumulates and interviews lengthen---and
saturates at finite certification cost plus the abstention penalty as physical
irrecoverability probability (dashed, right axis) approaches one. The illustrated
illustrated $\rho$ is sound exhaustive: it exhausts all declared paths and accepts
only evidence-established answers. Thus abstention and physical irrecoverability
coincide on this curve; that equality is a property of this policy, not a
definition. The shape---plateaus, cliffs, hazard-driven slope, saturation---is the
qualified prediction of \cref{prop:steps}; none of it is compound growth. This
worked curve sets $\Pr(\Err_q)=0$ to isolate decay; the simulation prices the
accepted-error channel separately.}
\label{fig:interest}
\end{figure}

\Cref{fig:interest} plots the resulting expected debt contribution and irrecoverability
probability against deficit age. Three features deserve narration because they are
the operational content of ``interest.'' First, \emph{the plateau}: for 90 days
the deficit is cheap to repay---this is the window in which proactive backfill
(\cref{sec:disc-management}) buys the most for the least. Second, \emph{the
cliffs}: cost multiplies at dates that are \emph{knowable in advance}---retention
expirations are published policy, migrations are scheduled projects---which is why
\cref{met:interest} replaced a fitted rate with the decay schedule itself: the
schedule is actionable (extend retention, export before migration), a rate is not.
Third, \emph{the saturation}: past the last non-human path, the sound exhaustive policy's
burden is dominated by its abstention penalty; further aging changes little because
the physical loss boundary has already been crossed (\cref{prop:writeoff}).

\subsection{Evidence Half-Life}
\label{sec:interest-halflife}

The replacement summary of \cref{met:interest} is computed the same way: given an
estate's path inventory and survival schedule, the \emph{evidence half-life}
$\tau_{1/2}$ is the age at which half of a reference query set ceases to be
answerable at baseline cost (first path death), and the \emph{irrecoverability
horizon} the age distribution of last path deaths. The worked example, having a
single shared schedule, makes $\tau_{1/2}$ degenerate to its one retention cliff
(90 days); the summary earns its ``half'' only over a heterogeneous estate whose
query set spans many schedules, which is its intended use. The structural point
survives the degeneracy: for the paradigmatic deficit, the entire first cliff is
set by a retention default that most organizations never chose
deliberately~\cite{nist80092}. That framing---reading governance-critical decay
dates off vendor defaults, rather than fitting a rate---is what makes the
schedule actionable: the cheapest evidence-debt intervention is often a retention
or export setting, not a new system.

\subsection{Hypotheses for Field Confirmation}
\label{sec:interest-hyps}

The interest analysis yields field-testable hypotheses, taken up in
\cref{sec:field}: \textbf{IH1}, measured reconstruction cost for comparable
queries has plateaus and cliffs aligned to the estate's
documented retention and migration schedule; \textbf{IH2}, reconstructions that
cross a personnel-departure boundary for the relevant service cost a multiple of
those that do not~\cite{rigby2016}; \textbf{IH3}, organizations' realized
irrecoverability concentrates in deficits older than the last surviving automated
path, with human memory as the modal final path. These are stated as hypotheses;
nothing in this paper confirms them.
